% ==============================================================================
% 1. Feladat: Rendszermodellezés és Paraméterillesztés
% ==============================================================================

\section{Feladat: Rendszermodellezés és Paraméterillesztés}

% Demonstration of the inline Lua computation bridge
\begin{luacode*}
function compute_parameters()
    local s_dist = 200.0   -- mm
    local phi_deg = 2.0    -- degree
    local phi_rad = phi_deg * math.pi / 180.0
    local f_target = s_dist * math.tan(phi_rad / 2.0)
    tex.print(string.format("\\def\\calcFocLength{%.2f}", f_target))
end
\end{luacode*}
\directlua{compute_parameters()}

\begin{problem}[15 pont]
Határozza meg a mérőrendszer bemeneti paramétereit a megadott \NeptunCode\ Neptun-kód alapján, majd vezesse le a finomállító pozicionáló tag átviteli függvényét az elektromechanikai differenciálegyenletekből!
\end{problem}

\subsection{Megadott feladatparaméterek}
A feladathoz tartozó kiindulási értékeket az \textbf{1. táblázat} foglalja össze a Style A (3 oszlopos, fehér rácsvonalas) specifikáció szerint \cite{bme_mogi_meresi_utmutato}.

% --- STYLE A: 3-Column Input Parameters Table ---
\begin{table}[ht!]
\centering
\renewcommand{\arraystretch}{1.25}
\setlength{\tabcolsep}{8pt}
\setlength{\arrayrulewidth}{2pt}
\arrayrulecolor{white}
\rowcolors{2}{white}{tableGray}
\begin{tabularx}{\textwidth}{|X|X|X|}
\hline
\rowcolor{white}
\textbf{Megnevezés} & \textbf{Jelölés} & \textbf{Adat} \\ \hline
Tárgy látószöge & $\varphi$ & $2^\circ$ \\ \hline
Legkisebb detektálandó méret & $x_{\min}$ & $0{,}50\text{ mm}$ \\ \hline
Tárgytávolság & $s$ & $200\text{ mm}$ \\ \hline
Laterális mozgási sebesség & $v$ & $0{,}01\text{ m/s}$ \\ \hline
Elméleti gyújtótávolság (Lua) & $f_{\text{calc}}$ & $\calcFocLength\text{ mm}$ \\ \hline
Megvilágítás típusa & --- & Természetes monokróm \\ \hline
\end{tabularx}
\caption{A megadott bemeneti paraméterek a \NeptunCode\ Neptun-kód alapján \cite{bme_mogi_meresi_utmutato}}
\label{tab:megadott_adatok}
\end{table}

\subsection{Matematikai levezetés és átviteli függvény}
A pozicionáló elektromechanikai alrendszer dinamikai egyenlete a Newton--Euler egyenlet és a huroktörvény alapján írható fel:
\begin{equation}
  J \ddot{\theta}(t) + b \dot{\theta}(t) + k \theta(t) = K_t \cdot i(t),
  \label{eq:mechanical_motion}
\end{equation}
ahol $J$ a redukált tehetetlenségi nyomaték, $b$ a csillapítási tényező, $k$ a visszatérítő rugómerevség, $K_t$ pedig a motor nyomatéktényezője.

Alkalmazva a Laplace-transzformációt zérus kezdeti feltételek mellett:
\begin{equation}
  (J s^2 + b s + k) \Theta(s) = K_t I(s).
\end{equation}

A rendszer szakasz-átviteli függvénye a szögelfordulás $\Theta(s)$ és a motoráram $I(s)$ között:
\begin{equation}
  \mathbf{P(s) = \frac{\Theta(s)}{I(s)} = \frac{\frac{K_t}{J}}{s^2 + \frac{b}{J}s + \frac{k}{J}} = \frac{K}{T_0^2 s^2 + 2\zeta T_0 s + 1}},
  \label{eq:transfer_function}
\end{equation}
ahol az időállandó $T_0 = \sqrt{J/k}$ és a csillapítási viszony $\zeta = \frac{b}{2\sqrt{J k}}$.
